% !TeX root = 0_TCC.tex
% Metodologia v2 — 01/11/2026. Não compilar isoladamente.
% Mantém o padrão do template e a bibliografia já instalada.
\begingroup
\justifying
\setlength{\parindent}{1.25cm}
\setlength{\parskip}{0pt}
\setlength{\emergencystretch}{2em}
\onehalfspacing
\renewcommand{\tablename}{Tabela}
\captionsetup[table]{labelsep=endash,justification=raggedright,singlelinecheck=false,font=small}

\phantomsection
\addcontentsline{toc}{section}{3 METODOLOGIA}
\noindent\textbf{3 METODOLOGIA}\par
\nobreak\vspace{\baselineskip}\nobreak

Esta pesquisa avalia se um modelo de visão computacional consegue localizar instrumentos cirúrgicos em fotografias e contar quantos exemplares de cada categoria estão presentes. Para isso, foi organizado um experimento com imagens da internet e do Hospital Universitário da Universidade Estadual de Londrina (HU-UEL). Cada fotografia constitui uma entrada para o modelo; cada instrumento marcado nessa fotografia constitui um objeto de referência. Essa distinção permite avaliar separadamente o reconhecimento dos instrumentos e o acerto das quantidades.

O procedimento foi organizado em cinco etapas: seleção das fontes; limpeza e anotação das imagens; reunião e separação dos dados; treinamento do modelo; e avaliação da detecção e da contagem. A sequência parte da preparação dos exemplos que o modelo utilizará para aprender e termina na análise de fotografias reservadas para o teste. Assim, o desempenho será examinado em dados que não participam do ajuste dos pesos da rede.

Os instrumentos selecionados têm como referência a caixa de apendicectomia, isto é, o conjunto de materiais utilizado na cirurgia de retirada do apêndice. A aplicação pretendida abrange a conferência de instrumentos no expurgo e no preparo, ambientes que integram o fluxo de processamento de materiais descrito por \textcite{moreto2025}. O experimento avalia a contagem em fotografias, sem abranger esterilidade, integridade física ou funcionamento dos instrumentos. Portanto, a viabilidade investigada é a de um apoio computacional à conferência, dentro das condições representadas pela base de imagens.

\vspace{\baselineskip}
\phantomsection
\addcontentsline{toc}{subsection}{3.1 Constituição e procedência da base de imagens}
\noindent\textbf{3.1 Constituição e procedência da base de imagens}\par
\nobreak\vspace{\baselineskip}\nobreak

A base de imagens reuniu duas origens complementares. Os conjuntos públicos ofereceram exemplos previamente anotados de diferentes instrumentos, enquanto as fotografias do HU-UEL acrescentaram registros dos materiais nos ambientes em que se pretende utilizar a contagem. As fontes da internet receberam os identificadores DS01 a DS07 para permitir acompanhar sua contribuição durante a organização. Esses identificadores são utilizados nas tabelas deste capítulo.

Os conjuntos DS01, DS03, DS04 e DS05 foram obtidos no Roboflow; o DS06, no Kaggle; e o DS07, no Hugging Face. Os endereços das fontes são indicados em nota.\footnote{DS01: \url{https://universe.roboflow.com/fmsurg/isef-training-dataset-p360u/dataset/1}; DS03: \url{https://universe.roboflow.com/surgical-tools-hnzri/surgical_tools-6b7ze}; DS04: \url{https://universe.roboflow.com/new-workspace-fjl7e/set_1_11/dataset/2}; DS05: \url{https://universe.roboflow.com/practicum-ii-iii/suturing-kit-tools}; DS06: \url{https://www.kaggle.com/datasets/dilavado/labeled-surgical-tools}; DS07: \url{https://huggingface.co/datasets/joonhaim/surgical-tool-recognition-full-multiview}.} O DS02 foi descartado integralmente após a revisão identificar imagens já presentes no DS01, inclusive versões com alterações de borda e ampliação. Dessa forma, seis conjuntos externos contribuíram para a base final.

As fotografias do HU-UEL foram organizadas inicialmente por instrumento e, em casos específicos, por composição de objetos. Entre as composições anotadas estão imagens com pinça dente de rato, pinça Backhaus e cúpula, além de imagens contendo uma tesoura Metzenbaum reta e outra curva. Também foram incorporadas 100 fotografias de bandejas vazias como exemplos negativos. Imagens sem instrumentos possuem arquivo de anotação vazio; por isso, a ausência de objetos nesses casos é intencional e não deve ser interpretada como falha de rotulagem.

As fotografias do HU-UEL foram produzidas por uma enfermeira em diferentes dias, ao longo de aproximadamente três semanas, nos ambientes de expurgo e preparo. Foi utilizado um \textbf{iPhone xx}. No primeiro dia, parte das imagens foi registrada em HEIC; posteriormente, o aparelho foi configurado no modo de alta compatibilidade, passando a produzir arquivos JPEG. As primeiras imagens foram aproveitadas após conversão para um formato compatível com o fluxo de anotação. Os registros utilizados mostram instrumentos e superfícies de apoio, sem pessoas ou procedimentos cirúrgicos em andamento.

A produção em dias e ambientes distintos acrescenta variação às fotografias, mas não permite assumir que todos os registros sejam independentes: imagens próximas podem mostrar o mesmo instrumento, sobre a mesma superfície, com pequenas mudanças de posição. Essa característica foi considerada na busca de repetições e na separação dos dados, descritas na Seção 3.3.

A Tabela~\ref{tab:met-origens} apresenta a composição após a seleção e a revisão. Essas quantidades descrevem os dados utilizados no experimento, e não resultados de desempenho do modelo. A contagem foi realizada por scripts Python que verificaram os pares de imagens e anotações. Uma fotografia com três instrumentos acrescenta uma unidade ao total de imagens e três ao total de objetos; uma bandeja vazia acrescenta apenas uma imagem.

\begin{table}[htbp]
\centering
\begin{minipage}{0.92\linewidth}
\singlespacing\small
\caption{Composição da base consolidada por origem}
\label{tab:met-origens}
\centering
\begin{tabular}{lrr}
\hline
Origem & Imagens & Objetos anotados \\
\hline
DS01 & 289 & 414 \\
DS03 & 115 & 129 \\
DS04 & 249 & 382 \\
DS05 & 85 & 85 \\
DS06 & 2.964 & 4.149 \\
DS07 & 1.008 & 1.008 \\
Subtotal da internet & 4.710 & 6.167 \\
HU-UEL & 5.617 & 6.547 \\
\hline
Total & 10.327 & 12.714 \\
\hline
\end{tabular}\par
\vspace{0.35em}\raggedright
Fonte: elaboração própria, a partir da revisão e contagem dos dados por scripts Python (2026).\par
Nota: as 100 imagens negativas estão incluídas no total do HU-UEL e não acrescentam objetos.
\end{minipage}
\end{table}

\vspace{\baselineskip}
\phantomsection
\addcontentsline{toc}{subsection}{3.2 Seleção, anotação e padronização dos instrumentos}
\noindent\textbf{3.2 Seleção, anotação e padronização dos instrumentos}\par
\nobreak\vspace{\baselineskip}\nobreak

Antes do treinamento, as imagens e suas anotações passaram por limpeza e conferência. Essa preparação é necessária porque o modelo aprende a partir dos exemplos disponíveis: uma classe incorreta ensina uma associação inadequada, um instrumento sem caixa pode ser interpretado como fundo e uma repetição entre treino e teste pode tornar a avaliação excessivamente favorável. A revisão procurou reduzir esses problemas antes de reunir os conjuntos.

Foram retiradas imagens de baixa qualidade que impediam identificar o instrumento, materiais incompatíveis com a caixa selecionada e versões repetidas reconhecidas durante a inspeção. Entre as repetições estavam imagens com espelhamento, alterações de perspectiva ou regiões cobertas por quadrados escuros, provenientes de aumento de dados aplicado previamente. Sempre que possível, foi preservada a versão original de melhor qualidade. A seleção considerou o conteúdo da fotografia, e não apenas seu nome ou o nome da classe no conjunto de origem. Dúvidas sobre instrumentos específicos foram esclarecidas com a enfermeira colaboradora.

Instrumentos visualmente semelhantes, mas fora do escopo, não foram incluídos apenas para aumentar o volume da base. Quando sua presença poderia gerar supervisão contraditória, optou-se pela retirada da imagem, inclusive se ela também apresentasse um objeto desejado. Um exemplo foi o afastador Army-Navy, cuja aparência em determinadas posições se aproxima daquela do Farabeuf. Essa escolha delimita o domínio do experimento: o detector não é avaliado como classificador universal de instrumentos semelhantes aos selecionados.

A revisão dos conjuntos externos foi realizada com apoio do CVAT e de inspeções no Roboflow. As fotografias do HU-UEL foram anotadas no X-AnyLabeling. Anotar uma imagem, neste estudo, significa traçar um retângulo que envolva cada instrumento e atribuir a ele uma classe. Esses retângulos, ou caixas delimitadoras, constituem as referências utilizadas para ensinar e avaliar o detector. Foram utilizadas caixas horizontais, com lados paralelos às bordas da fotografia, mesmo quando o instrumento estava inclinado.

As anotações foram exportadas no formato YOLO para detecção. Cada linha do arquivo TXT representa um objeto e contém o identificador da classe, as coordenadas do centro da caixa, sua largura e sua altura. As quatro medidas são divididas pelas dimensões da imagem, ficando normalizadas entre zero e um. Por exemplo, uma largura de 0,25 representa um quarto da largura da fotografia. A imagem e seu TXT compartilham o mesmo nome-base, permitindo associá-los sem depender da ordem dos arquivos na pasta.

Nos conjuntos com polígonos, a conversão para detecção utilizou os menores e maiores valores de cada coordenada para obter a caixa envolvente de cada instância. Esse procedimento conserva a extensão espacial do contorno anotado, mas não verifica sua pertinência semântica. Por isso, caixas que abrangiam componentes indesejados, como tubos conectados ao bico de aspiração ou partes além do cabo de bisturi, foram submetidas à revisão manual. No DS07, as anotações convertidas de JSONL foram inspecionadas visualmente, e foram mantidas as classes correspondentes a porta-agulhas e cabos de bisturi.

Para utilizar os mesmos dados em diferentes níveis de detalhe, foram preservadas duas organizações das classes, também chamadas de taxonomias. A específica distingue 17 instrumentos ou variantes, incluindo Kelly, Mayo e Metzenbaum retas e curvas. A geral reúne essas classes em sete categorias e foi utilizada no treinamento principal. A correspondência é fixa: cada objeto de uma classe específica recebe uma única categoria geral. O agrupamento reduz o número de nomes que o modelo precisa distinguir, mas conserva diferenças visuais dentro de categorias como pinças; não se pressupõe que todos os instrumentos agrupados tenham a mesma aparência.

Diferenças de tamanho não originam classes distintas. Porta-agulhas, cúpula e bico de aspiração foram mantidos em categorias próprias. A pinça Collin coração não integra esta etapa por falta de imagens, e a Kocher curva não pertence ao conjunto selecionado. Os exemplares Yankauer aceitos na conferência foram reunidos à categoria de bicos de aspiração.

A Tabela~\ref{tab:met-taxonomias} explicita a correspondência entre as taxonomias e a quantidade de objetos em cada categoria geral. Os nomes em português são utilizados na apresentação do trabalho; os identificadores em inglês são empregados nos arquivos de treinamento. As duas organizações mantêm as mesmas caixas e divisões, alterando apenas os identificadores de classe. Portanto, não constituem duas bases independentes.

\begin{table}[htbp]
\centering
\begin{minipage}{\linewidth}
\singlespacing\small
\caption{Correspondência entre classes específicas e categorias gerais}
\label{tab:met-taxonomias}
\begin{tabular}{p{0.19\linewidth}p{0.61\linewidth}r}
\hline
Categoria geral & Classes específicas representadas & Objetos \\
\hline
Pinças & Allis; anatômica; Backhaus; Kelly curva; Kelly reta; Foerster; dente de rato; Kocher reta & 4.907 \\
Tesouras & Mayo curva; Mayo reta; Metzenbaum curva; Metzenbaum reta & 3.976 \\
Afastadores & Farabeuf & 848 \\
Porta-agulhas & Porta-agulhas & 760 \\
Cabos de bisturi & Cabo de bisturi & 1.673 \\
Bicos de aspiração & Bico de aspiração, incluindo exemplares Yankauer aceitos na revisão & 244 \\
Cúpulas & Cúpula & 306 \\
\hline
Total & 17 classes específicas e 7 categorias gerais & 12.714 \\
\hline
\end{tabular}\par
\vspace{0.35em}\raggedright
Fonte: elaboração própria, a partir da revisão e contagem dos dados por scripts Python (2026).
\end{minipage}
\end{table}

A representação é desigual entre classes e origens. Por exemplo, as pinças anatômicas e os cabos de bisturi da base consolidada provêm exclusivamente da internet. Consequentemente, um resultado agregado satisfatório não demonstra desempenho hospitalar para todas as classes. Essa diferença deve ser considerada na análise dos erros e na apresentação dos resultados por origem.

\vspace{\baselineskip}
\phantomsection
\addcontentsline{toc}{subsection}{3.3 Consolidação, verificação e separação dos dados}
\noindent\textbf{3.3 Consolidação, verificação e separação dos dados}\par
\nobreak\vspace{\baselineskip}\nobreak

Neste estudo, consolidação significa reunir os conjuntos selecionados em uma única base organizada, compatibilizando nomes de classes, formatos de anotação e estrutura de pastas. Esse processo foi executado por scripts Python sobre cópias dos arquivos, preservando os conjuntos originais. Cada imagem consolidada manteve uma identificação de sua origem, de suas classes e da divisão à qual foi destinada. A preparação ocorreu no Windows, e a base organizada foi transferida para o Mac utilizado no treinamento.

A conferência automática verificou se cada imagem podia ser aberta, se possuía a anotação correspondente e se os identificadores de classe estavam previstos no mapeamento. Também foram examinadas as coordenadas das caixas, incluindo dimensões positivas e limites válidos. Os arquivos vazios foram tratados separadamente para distinguir as bandejas sem instrumentos de uma possível ausência de anotação. As duas taxonomias foram comparadas para verificar a manutenção das mesmas imagens, caixas e divisões, com alteração apenas das classes.

A busca por repetições combinou comparação de arquivos, comparação dos pixels decodificados e semelhança visual. Para isso, foram utilizados resumos digitais, chamados \textit{hashes}: o SHA-256 permitiu conferir a integridade dos arquivos comparados, enquanto o dHash forneceu uma assinatura simplificada da aparência. Assinaturas perceptuais próximas indicaram candidatos à revisão, mas não foram tratadas como prova suficiente de duplicação. Fotografias de instrumentos diferentes sobre o mesmo fundo podem apresentar grande semelhança nessa representação.

Os candidatos foram examinados visualmente para distinguir cópias de fotografias diferentes. Quando uma mesma imagem apareceu em resoluções distintas, a versão de melhor qualidade foi priorizada. Fotografias relacionadas por repetição ou sequência identificada foram reunidas em grupos destinados à mesma divisão. Dessa forma, procurou-se evitar que uma versão fosse utilizada no aprendizado e outra muito próxima fosse usada para avaliar o mesmo modelo.



O agrupamento procura reduzir a dependência entre dados usados para aprender e dados usados para avaliar. Como discutem \textcite{kapoor2023}, falhas na separação podem produzir conclusões excessivamente otimistas. Entretanto, a ausência de grupos cadastrados em mais de uma divisão não comprova independência de todas as sessões de aquisição. A base contém imagens de instrumentos reposicionados em cenários semelhantes, e os registros disponíveis não permitem reconstruir integralmente todas as sessões das fontes externas. Essa restrição permanece como limitação do experimento.

Após a organização, a base foi separada em treino, validação e teste. O treino fornece os exemplos utilizados para ajustar os pesos da rede. A validação acompanha o aprendizado e permite selecionar o estado do modelo e o limiar de confiança. O teste permanece reservado para avaliar o procedimento após essas escolhas. A distribuição final se aproxima de 80\%, 10\% e 10\%, respectivamente, com pequenos ajustes necessários para manter os grupos identificados na mesma divisão.

A Tabela~\ref{tab:met-splits} apresenta as quantidades efetivas. A mesma divisão foi preservada nas duas taxonomias, permitindo uma eventual comparação sobre as mesmas fotografias. Entre os exemplos negativos, 80 bandejas vazias pertencem ao treino, dez à validação e dez ao teste. A separação por grupos reduz um risco de dependência, mas não equivale a uma divisão completa por dia, paciente ou hospital: as fontes disponíveis não fornecem todos esses metadados.

\begin{table}[htbp]
\centering
\begin{minipage}{0.92\linewidth}
\singlespacing\small
\caption{Distribuição das imagens e dos objetos por divisão}
\label{tab:met-splits}
\centering
\begin{tabular}{lrr}
\hline
Divisão & Imagens & Objetos anotados \\
\hline
Treino & 8.304 & 10.044 \\
Validação & 1.015 & 1.293 \\
Teste & 1.008 & 1.377 \\
\hline
Total & 10.327 & 12.714 \\
\hline
\end{tabular}\par
\vspace{0.35em}\raggedright
Fonte: elaboração própria, a partir da revisão e contagem dos dados por scripts Python (2026).
\end{minipage}
\end{table}

\vspace{\baselineskip}
\phantomsection
\addcontentsline{toc}{subsection}{3.4 Ambiente computacional e configuração do treinamento}
\noindent\textbf{3.4 Ambiente computacional e configuração do treinamento}\par
\nobreak\vspace{\baselineskip}\nobreak

O ambiente de execução utiliza um Mac mini com chip Apple M4 e 16~GB de memória unificada, com aceleração pela interface Metal Performance Shaders (MPS) do PyTorch. O ambiente utilizado compreende Python 3.13, PyTorch 2.15.0 e Ultralytics 8.4.192, em arquitetura arm64. A configuração foi mantida após os testes de instalação para evitar mudanças de dependências durante o experimento.

Antes do treinamento principal, foi executado um piloto de três épocas com 10\% das imagens de treino. Sua finalidade foi verificar o carregamento das imagens e dos rótulos, a execução na GPU por MPS, o comportamento do consumo de memória e o tempo necessário para processar os lotes. Essa etapa permitiu verificar a execução do procedimento no equipamento disponível antes de iniciar uma tarefa mais longa.

O treinamento principal foi iniciado com o YOLO26n e pesos pré-treinados, utilizando toda a divisão de treino. Portanto, o aprendizado principal começou a partir do pré-treinamento, e não dos pesos obtidos no piloto. Como apresentado na fundamentação, a letra \textit{n} identifica a variante de menor porte da família utilizada, que também inclui variantes maiores. Foi priorizada uma configuração compatível com a memória e o tempo disponíveis, sem supor que ela apresentaria o melhor desempenho entre todas as alternativas.

O pré-treinamento fornece uma inicialização que já aprendeu características visuais em outra base. Durante a adaptação aos instrumentos cirúrgicos, os pesos podem ser atualizados, pois não foi configurado congelamento de camadas. Essa escolha reduz a necessidade de iniciar o aprendizado com parâmetros aleatórios, mas não substitui a presença de exemplos adequados das classes e dos cenários de interesse \parencite{pytorchTransferencia}.

A configuração principal está sintetizada na Tabela~\ref{tab:met-treino}. O limite de 50 épocas foi definido como orçamento inicial de treinamento, e não como valor ótimo estabelecido por busca sistemática. A parada antecipada foi configurada com tolerância de dez épocas sem melhora no critério de validação da implementação. A quantidade efetivamente executada, a época selecionada e a duração serão informadas a partir dos registros, após o encerramento.

\begin{table}[htbp]
\centering
\begin{minipage}{\linewidth}
\singlespacing\small
\caption{Parâmetros definidos para o treinamento principal}
\label{tab:met-treino}
\begin{tabular}{p{0.52\linewidth}p{0.40\linewidth}}
\hline
Parâmetro & Configuração \\
\hline
Modelo e tarefa & YOLO26n; detecção de caixas horizontais \\
Classes e fração do treino & 7 categorias; 100\% da divisão \\
Épocas e tolerância de parada & Máximo de 50; tolerância de 10 \\
Dimensão de entrada configurada & 640 pixels \\
Lote físico / lote nominal & 4 / 64 \\
Otimizador & AdamW \\
Taxa inicial / fator final & 0,001 / 0,01 \\
Decaimento de pesos configurado & 0,0005 \\
Parâmetro de momento & 0,9 \\
Aquecimento & 3 épocas \\
Semente & 42 \\
Dispositivo e precisão mista & MPS; precisão mista desativada \\
Processos auxiliares / cache de imagens & 0 / desativado \\
Salvamento periódico & A cada 5 épocas \\
\hline
\end{tabular}\par
\vspace{0.35em}\raggedright
Fonte: elaboração própria, a partir da configuração do treinamento (2026).\par
Nota: o lote nominal é uma referência interna para a otimização; não equivale ao número de imagens simultaneamente processadas no dispositivo.
\end{minipage}
\end{table}

O AdamW emprega decaimento de pesos desacoplado da atualização adaptativa, conforme a formulação de \textcite{loshchilov2019adamw}. Sua adoção neste experimento representa uma escolha explícita de otimização, sem reproduzir integralmente a receita de pré-treinamento do YOLO26. A taxa de aprendizagem foi configurada com aquecimento e redução linear, sem busca automática de hiperparâmetros. As definições operacionais dos argumentos são apresentadas na documentação do software \parencite{ultralyticsTreinamento}.

O tamanho de entrada foi configurado em 640 pixels para limitar o custo de processamento e utilizar uma dimensão usual da implementação. Durante o carregamento, as imagens são adaptadas às entradas do modelo por redimensionamento e, conforme a transformação aplicada, preenchimento ou recorte. Os arquivos de origem não são alterados por esse processamento. Não foi ativado treinamento com tamanhos variáveis entre lotes. A redução de resolução pode remover detalhes de pontas, serrilhas e instrumentos pequenos; por isso, o tamanho escolhido constitui uma condição do experimento, e não uma resolução comprovadamente ideal.

O lote físico foi limitado a quatro imagens, isto é, quatro entradas processadas de cada vez. O lote nominal de 64 é uma referência interna da implementação para a otimização e não indica que 64 imagens estejam simultaneamente carregadas na GPU. O cache de imagens e os processos auxiliares de carregamento foram desativados, assim como a precisão mista. Essas decisões priorizam uma configuração de execução compatível com o Mac utilizado; não foram comparadas sistematicamente com outras combinações.

A semente foi fixada em 42 e foi configurada execução determinística. A semente controla as sequências de escolhas aleatórias, contribuindo para documentar o experimento. Entretanto, algumas operações do MPS não dispõem de implementação determinística, e o próprio ambiente apresentou aviso dessa limitação. Assim, a configuração favorece a repetibilidade, mas não garante resultados numericamente idênticos em todas as execuções \parencite{pytorchReprodutibilidade}.

Durante o treino, a validação é realizada ao final das épocas. São registrados os argumentos efetivos, as versões dos programas, as perdas e as métricas de validação, além dos pesos do último estado e do estado considerado melhor pelo critério da implementação. O conjunto de teste não participa dessa seleção. Os registros permitem diferenciar a configuração planejada da quantidade de épocas e da duração efetivamente observadas.

\vspace{\baselineskip}
\phantomsection
\addcontentsline{toc}{subsection}{3.5 Aumento de dados durante o treinamento}
\noindent\textbf{3.5 Aumento de dados durante o treinamento}\par
\nobreak\vspace{\baselineskip}\nobreak

O aumento de dados cria variações das imagens durante o aprendizado, procurando reduzir a dependência de posições e aparências muito específicas \parencite{goodfellow2016}. Neste experimento, essas transformações foram configuradas para ocorrer durante o treino, sem salvar uma nova coleção permanente de fotografias. Assim, a mesma imagem pode ser apresentada de formas diferentes em épocas distintas, com atualização correspondente das caixas.

Foram configurados espelhamentos horizontal e vertical com probabilidade de 0,5 cada, rotações de até aproximadamente $\pm15^{\circ}$, translação relativa de até 0,1 e variação de escala de 0,25. Essas operações procuram representar mudanças de orientação, posição e tamanho aparente. As variações de matiz, saturação e brilho, controladas pelos parâmetros HSV, foram definidas como 0,015, 0,3 e 0,2. Elas introduzem diferenças de cor e iluminação, sem demonstrar, por si só, que o modelo será robusto a todas as condições do hospital.

A transformação mosaic foi configurada com probabilidade de 0,5 para compor entradas a partir de diferentes imagens. Sua aplicação busca ampliar a variedade de composições, mas não reproduz necessariamente oclusões e interações físicas entre instrumentos em uma bandeja real. MixUp, Copy-Paste, cisalhamento e perspectiva adicional foram desativados. A documentação da Ultralytics descreve essas operações e seus parâmetros \parencite{ultralyticsAugmentation}.

Foi prevista a desativação do mosaic nas dez últimas épocas da agenda programada. Caso a parada antecipada ocorra antes dessa etapa, não se pode afirmar que o modelo tenha passado por esse período final sem mosaic. As transformações efetivamente carregadas, inclusive componentes opcionais dependentes do ambiente, são registradas para conferir a correspondência com a configuração solicitada. Exemplos de lotes transformados permitem inspecionar cortes e caixas produzidos durante o treinamento.

As 8.304 imagens de treino permanecem como quantidade física de referência. Repetições ao longo das épocas e variações geradas durante o processamento não são contabilizadas como fotografias independentes. Validação e teste não recebem a política aleatória de aumento de dados do treino, embora também necessitem do pré-processamento de entrada. Essa distinção evita confundir expansão de apresentações do mesmo material com aumento da diversidade original da amostra.

\vspace{\baselineskip}
\phantomsection
\addcontentsline{toc}{subsection}{3.6 Protocolo de avaliação da detecção e da contagem}
\noindent\textbf{3.6 Protocolo de avaliação da detecção e da contagem}\par
\nobreak\vspace{\baselineskip}\nobreak

A avaliação final será realizada após o encerramento do treinamento e terá duas perguntas complementares: o modelo localiza e classifica os instrumentos corretamente? As detecções aceitas produzem a quantidade correta por imagem e por categoria? Essas perguntas precisam ser analisadas em conjunto, pois uma quantidade correta pode resultar da compensação entre uma omissão e uma detecção indevida.

O estado do modelo será selecionado a partir da validação. Nessa mesma divisão será definido o limiar de confiança, valor mínimo exigido para aceitar uma predição na contagem. O teste não será utilizado para procurar esse valor. Antes de sua execução, serão fixados os pesos, o tamanho de entrada, o limiar e as opções de pós-processamento. Essa separação evita ajustar o procedimento às fotografias utilizadas para apresentar seu desempenho final.
% PENDENCIA AVAL-01: registrar a regra de calibração e o limiar efetivamente adotados antes do teste; proposta no roteiro complementar, ainda não executada.

Para detecção, serão obtidos precisão, revocação, mAP a 0,50 de interseção sobre união e mAP de 0,50 a 0,95, com resultados por categoria. A interseção sobre união mede quanto a caixa prevista coincide com a caixa de referência. As métricas serão interpretadas conforme as definições da fundamentação e as diferenças discutidas por \textcite{padilla2021}. O cálculo padrão da implementação será registrado separadamente das medidas produzidas no limiar escolhido para a contagem, pois os procedimentos podem utilizar condições distintas \parencite{ultralyticsValidacao}.

As métricas de contagem serão calculadas por um script que leia as anotações de referência e as predições. Para cada imagem, o programa contará os objetos de cada categoria e comparará as quantidades. Serão obtidos o erro absoluto médio do total, a proporção de imagens com total correto, o erro por categoria e a proporção com acerto simultâneo em todas as categorias. Esta última medida compara o vetor de contagens, isto é, a lista de quantidades na ordem fixa das classes. Uma troca entre duas categorias pode preservar o total e, ainda assim, tornar esse vetor incorreto.

Também serão separadas as imagens com contagem inferior e superior à referência. Fotografias sem instrumentos serão incluídas, pois permitem verificar se o modelo acrescenta objetos indevidamente. As medidas serão apresentadas para o conjunto de teste e, quando houver dados suficientes, por origem e categoria. Serão informados os respectivos números de imagens e objetos, evitando interpretar a ausência de exemplos de uma classe como desempenho nulo ou perfeito.

O número de objetos por fotografia será obtido diretamente das anotações. A análise por quantidade distinguirá cenas vazias, com um objeto e com múltiplos objetos, permitindo observar se o comportamento muda conforme a composição. Já a sobreposição física, os reflexos, o desfoque e os cortes de instrumentos dependem de inspeção visual. A interseção entre caixas não será utilizada como prova de oclusão, pois caixas de objetos separados também podem se sobrepor.

A revisão qualitativa examinará exemplos de omissões, detecções excedentes, trocas de categoria e localização inadequada. As possíveis condições associadas às falhas serão descritas sem atribuir causalidade apenas pela inspeção. Caso seja realizada comparação quantitativa entre imagens com e sem oclusão, essa característica precisará ser anotada previamente nas imagens avaliadas. Sem essa classificação, a discussão sobre oclusões permanecerá qualitativa.

O custo computacional será descrito pela duração efetiva do treinamento e por medições de processamento no Mac utilizado. Na avaliação de tempo por imagem, será necessário distinguir a inferência da rede do tempo total, que também pode incluir leitura, preparação da imagem e organização das saídas. As primeiras execuções de aquecimento serão separadas das medições repetidas, e a sincronização do dispositivo será considerada para não medir apenas o envio de operações à GPU.

Um eventual experimento com as 17 classes específicas utilizará as mesmas divisões de imagens. Seus resultados de detecção serão apresentados como uma tarefa de maior detalhamento, sem tratar os valores de mAP das duas taxonomias como diretamente equivalentes. Para comparar contagens nas sete categorias, as predições específicas poderão ser reunidas pelo mesmo mapeamento da Tabela~\ref{tab:met-taxonomias}. A análise só incluirá execuções efetivamente concluídas.

A viabilidade será discutida a partir da combinação entre localização e classificação, acerto das contagens, frequência e natureza das falhas e custo de processamento. Não foi definido um limite clínico de aceitação. Portanto, a conclusão deverá explicitar as condições em que a contagem apresentou resultados favoráveis e aquelas em que permaneceu insuficiente, sem converter o desempenho em fotografias em autorização para substituir a conferência profissional. Novas fotografias de bandejas completas poderão ampliar a avaliação externa, caso estejam disponíveis, mas não serão tratadas como parte de um teste já executado.

\endgroup
